\section{边界的消弭：Browser、Desktop、Mobile、GUI、CLI、API 与 Coding}

苏煜反复强调，Agent 领域早期会区分 browser use、desktop use、mobile use、GUI、text-based representation、CLI、API、coding 等，但这些划分是临时性的。最终大家想要的是 universal digital agent。这些接口都是手段，而不是最终目的。

\begin{figure}[H]
\centering
\includegraphics[width=0.92\textwidth]{boundary-convergence.png}
\caption{边界的消弭：多种数字环境入口正在收敛到 universal digital agent。}
\end{figure}

\begin{knowledgebox}{读图：为什么 coding 是边界消融的核心}
图中所有入口都连向 universal digital agent。Browser、desktop、mobile 和 GUI 是表层交互；CLI、API、coding 是更形式化的控制接口。苏煜认为 coding 是 digital world 的 fabric，因为很多界面最终都可以被代码表达、渲染或操纵。
\end{knowledgebox}

\subsection{GUI 不会消失，CLI 也不会统治一切}

访谈没有走向“GUI 会被 CLI 全面取代”的极端。原因有两点。第一，现实世界有大量 legacy system，例如银行、企业软件和老系统，不会快速重写。第二，即使 CLI 对 Agent 是全局最优，也不意味着对所有局部场景都是最优。很多现有 GUI 对人类和组织来说已经 good enough。

\begin{warningbox}{局部最优与全局最优不同}
技术上更适合 Agent 的接口，不一定会立刻替代已有界面。企业迁移成本、用户习惯、监管、遗留系统和经济账都会决定实际路径。Agent 技术判断必须和部署环境一起看。
\end{warningbox}

\subsection{Programming Language 也是 Language}

主持人提出“自然语言是人类脚手架，coding 是机器脚手架”的比喻。苏煜进一步指出，language 从来不只是自然语言，编程语言、图表、手势都是广义 language。Programming language 是 formal language，它更精确、更可执行，因此在 Agent 操作 digital world 时尤其重要。

\subsection{本章小结}

Agent 的终局不是 browser agent、desktop agent 或 coding agent 的单项胜利，而是多种入口在任务层收敛。Coding 的重要性在于它把许多接口统一成可表达、可组合、可执行的对象。

